\scanpage{50}
\figure{\hfill\tikz[x=.52cm,y=.52cm]{\draw[thin] (0,0)--(-.12,1.6)--(.43,.03)--(.31,1.84)--(.82,.17) .. controls (.88,-.04) and (.8,-.12) .. (.77,-.13); \draw[thin] (.83,.67)--(.97,1.34) (1.04,.61)--(1.22,1.41);}}
\mspage{- 82 -}\hfill n\textsuperscript{o}357

\[
\deleted{f^\rho(x\sigma)=\chi(\sigma)^{-1}f^{\handwritten{\rho}}(x)}
\]
\deleted{\uline{quel que soit} $\sigma\in H$.}

\deleted{Soit $K$ compact dans $X$ ; alors $x\in K$, $x\xi\in S$ entraîne $x\xi\in S\cap KH$, d'où, puisque $S\cap KH$ est compact et que $H$ opère proprement, $\xi\in L$, où $L$ est une partie compacte de $H$ qui ne dépend que de $K$. Pour $x\in K$, on peut donc remplacer $\int_H$ par $\int_L$ dans (1) ; d'après le lemme 2 du Chap.III, \S\,5, n\textsuperscript{o}1, il s'ensuit que $f^{\handwritten{\rho}}$ est continue sur $K$. Le reste est évident.}

\deleted{Si une fonction $g$ satisfait à $g(x\sigma)=\chi(\sigma)^{-1}g(x)$ quels que soient $x\in X$, $\sigma\in H$, son support $S$ est invariant par $H$, et s'écrit donc $\pi^{-1}(\pi(S))$. On désignera par $\handwritten{\unsure{\mathscr K|\mathcal K}}^{\handwritten{\chi}}(X,H)$ l'ensemble des fonctions continues $g$ sur $X$, à valeurs dans $\underline{\underline{\mathrm R}}$, qui satisfont à $g(x\sigma)=\chi(\sigma)^{-1}g(x)$ quels que soient $x\in X$ et $\sigma\in H$, et dont le support est de la forme $KH$ avec $K$ compact ; il est clair que c'est un espace de Riesz ; pour $\chi=1$, $h\to h\circ\pi$ est un isomorphisme de $\handwritten{\mathscr K}(X/G)$ sur $\handwritten{\mathscr K}^{1}(X,G)$. La prop.1 montre que $f\to f^\rho$ détermine une application de $\handwritten{\mathscr K}(X)$ dans $\handwritten{\mathscr K}^{\handwritten{\chi}}(X,G)$, qui est manifestement \uline{croissante}. On désigne par $\handwritten{\mathscr K}_+^{\handwritten{\chi}}(X,H)$ l'ensemble des fonctions $\geq 0$ appartenant à $\handwritten{\mathscr K}^{\handwritten{\chi}}(X,H)$.}

\deleted{PROPOSITION 2.- \uline{Soit} $K$ \uline{une partie compacte de} $X$. \uline{Alors il existe une fonction} $h\in\handwritten{\mathscr K}_+(X)$ \uline{telle que}, \uline{pour toute fonction} $g\in\handwritten{\mathscr K}^{\handwritten{\chi}}(X,H)$ \uline{de support contenu dans} $KH$, \uline{on ait} $g=(hg)^{\handwritten{\rho}}$.}

\deleted{Soit $f\in\handwritten{\mathscr K}_+(X)$ telle que $f>0$ sur $K$ ; posons :}
\[
\deleted{f'(x)=\int_H f(x\xi)\chi(\xi)^{-1}\mathrm d\handwritten{\rho}(\xi)\ ;}
\]
\deleted{comme $\chi^{-1}\handwritten{\rho}$ est une mesure de Haar sur $H$, $f'$ est continue, d'après le cor. de la prop.1, et de la forme $f^b\circ\pi$, avec $f^b\in\handwritten{\mathscr K}(X/H)$. Soit $U$ l'ensemble ouvert des points où $f^b>0$ ; comme $f'>0$ sur $K$, on a $U\supset\pi(K)$ ; $f/f'$ est continue sur $\pi^{-1}(U)$. Soit $U'$ une partie ouverte de $X/H$, telle que $\pi(K)\subset U'$ et $\overline U'\subset U$ ; soit $F$ continue $\geq 0$ sur $X/H$, égale à 1 sur $\pi(K)$ et}
% archive sheet number: 49
